\documentclass[11pt]{ctexart}
% cnreport-preamble.tex — 中文技术报告导言区(Cosmos/arXiv 单栏风格)
% 用法: \documentclass[11pt]{ctexart} 之后 % cnreport-preamble.tex — 中文技术报告导言区(Cosmos/arXiv 单栏风格)
% 用法: \documentclass[11pt]{ctexart} 之后 % cnreport-preamble.tex — 中文技术报告导言区(Cosmos/arXiv 单栏风格)
% 用法: \documentclass[11pt]{ctexart} 之后 \input{cnreport-preamble}
% 编译: tectonic(XeTeX 引擎),需要系统安装 Noto Serif/Sans CJK SC 字体
\usepackage[a4paper,top=2.7cm,bottom=2.7cm,left=2.5cm,right=2.5cm,headheight=15pt]{geometry}
\usepackage{fontspec}
% 正文:拉丁衬线 + 思源宋体;标题:拉丁无衬线 + 思源黑体
% 按文件名从 TeX 字体树加载(tectonic bundle 自带 TeX Gyre;系统字体需 Noto CJK SC)
\setmainfont{texgyretermes}[Extension=.otf,UprightFont=*-regular,BoldFont=*-bold,ItalicFont=*-italic,BoldItalicFont=*-bolditalic]
\setsansfont{texgyreheros}[Extension=.otf,UprightFont=*-regular,BoldFont=*-bold,ItalicFont=*-italic,BoldItalicFont=*-bolditalic]
\setmonofont{texgyrecursor}[Extension=.otf,UprightFont=*-regular,BoldFont=*-bold,Scale=0.88]
\setCJKmainfont{Noto Serif CJK SC}[BoldFont={Noto Serif CJK SC Bold}]
\setCJKsansfont{Noto Sans CJK SC}[BoldFont={Noto Sans CJK SC Bold}]
\setCJKmonofont{Noto Sans Mono CJK SC}
% 带圈数字 ①-⑳ 等划入 CJK 类,用 CJK 字体渲染(xeCJK 默认按拉丁处理)
\xeCJKDeclareCharClass{CJK}{"2460 -> "24FF}

\usepackage{amsmath,amssymb,bm}
\usepackage{booktabs}
\usepackage{graphicx}
\usepackage{xcolor}
\usepackage{titlesec}
\usepackage{fancyhdr}
\usepackage{enumitem}
\usepackage{caption}
\usepackage{listings}

% ---- 配色(主色可改) ----
\definecolor{accent}{HTML}{0E5FA8}     % 强调蓝(链接/节号)
\definecolor{good}{HTML}{0E7A4D}       % 正向结果
\definecolor{warn}{HTML}{B45309}       % 负向/警示
\definecolor{faint}{HTML}{6B7889}      % 弱化文字
\definecolor{codebg}{HTML}{F0F2F5}
\definecolor{rule}{HTML}{C9CFD8}

\usepackage[colorlinks=true,linkcolor=accent,citecolor=accent,urlcolor=accent]{hyperref}

% ---- 章节标题:无衬线粗体,节号着色 ----
\titleformat{\section}{\Large\bfseries\sffamily}{\textcolor{accent}{\thesection}}{0.6em}{}
\titleformat{\subsection}{\large\bfseries\sffamily}{\textcolor{accent}{\thesubsection}}{0.5em}{}
\titleformat{\subsubsection}{\normalsize\bfseries\sffamily}{\textcolor{accent}{\thesubsubsection}}{0.5em}{}
\titlespacing{\section}{0pt}{1.6em}{0.7em}
\titlespacing{\subsection}{0pt}{1.2em}{0.5em}

% ---- 页眉页脚(Cosmos 式:左标题右日期,下细线) ----
\newcommand{\runningtitle}{}  % 主文档用 \renewcommand 设置
\newcommand{\reportdate}{}
\newcommand{\reportfooter}{}
\fancypagestyle{cnreport}{%
  \fancyhf{}
  \fancyhead[L]{\small\sffamily \runningtitle}
  \fancyhead[R]{\small\sffamily \reportdate}
  \fancyfoot[L]{\small\color{faint}\sffamily \reportfooter}
  \fancyfoot[R]{\small\sffamily \thepage}
  \renewcommand{\headrulewidth}{0.4pt}
  \renewcommand{\footrulewidth}{0pt}
}
\pagestyle{cnreport}

% ---- 题注与列表 ----
\captionsetup{font=small,labelfont={bf,sf},skip=6pt}
\setlist{itemsep=2pt,topsep=4pt,parsep=0pt}

% ---- 代码 ----
\lstset{
  basicstyle=\ttfamily\small,
  backgroundcolor=\color{codebg},
  frame=single, framerule=0.3pt, rulecolor=\color{rule},
  xleftmargin=6pt, xrightmargin=6pt,
  breaklines=true, columns=fullflexible, keepspaces=true,
  aboveskip=8pt, belowskip=8pt,
}

% ---- 报告头(kicker / 标题 / 副题 / 署名 / 摘要 / 关键词) ----
% 用法: \reporthead{kicker-left}{kicker-right}{标题}{副题}{署名行}
\newcommand{\reporthead}[5]{%
  \begin{center}
  {\small\sffamily\bfseries\color{accent} #1 \hfill #2}\\[2pt]
  {\color{black}\rule{\linewidth}{1.2pt}}\\[14pt]
  {\LARGE\sffamily\bfseries #3\par}\vspace{8pt}
  {\large\color{faint} #4\par}\vspace{10pt}
  {\small\sffamily\color{faint} #5\par}\vspace{4pt}
  {\color{rule}\rule{\linewidth}{0.4pt}}
  \end{center}\vspace{6pt}
}
% 摘要块: \begin{cnabstract} ... \end{cnabstract}
\newenvironment{cnabstract}{%
  \begin{quote}\small\color{faint}\sffamily\bfseries\color{black} 摘要 \quad\normalfont\rmfamily\color{faint!80!black}%
}{\end{quote}\vspace{2pt}}
% 重点结论框: \begin{quotebox} ... \end{quotebox}
\newenvironment{quotebox}{%
  \begin{quote}\small
  {\color{rule}\rule{\linewidth}{0.8pt}}\par\vspace{2pt}%
}{\par\vspace{2pt}{\color{rule}\rule{\linewidth}{0.8pt}}\end{quote}\vspace{2pt}}

\newcommand{\keywords}[1]{\begin{quote}\small\color{faint}{\bfseries\sffamily\color{black} 关键词:\ } #1\end{quote}}

% 表内彩色标记
\newcommand{\upgood}[1]{\textcolor{good}{\bfseries #1}}
\newcommand{\downbad}[1]{\textcolor{warn}{\bfseries #1}}
\newcommand{\dimtext}[1]{\textcolor{faint}{#1}}

\setlength{\parskip}{2pt}
\linespread{1.12}
\emergencystretch=3em  % CJK 段落容忍度,消除 overfull

% 编译: tectonic(XeTeX 引擎),需要系统安装 Noto Serif/Sans CJK SC 字体
\usepackage[a4paper,top=2.7cm,bottom=2.7cm,left=2.5cm,right=2.5cm,headheight=15pt]{geometry}
\usepackage{fontspec}
% 正文:拉丁衬线 + 思源宋体;标题:拉丁无衬线 + 思源黑体
% 按文件名从 TeX 字体树加载(tectonic bundle 自带 TeX Gyre;系统字体需 Noto CJK SC)
\setmainfont{texgyretermes}[Extension=.otf,UprightFont=*-regular,BoldFont=*-bold,ItalicFont=*-italic,BoldItalicFont=*-bolditalic]
\setsansfont{texgyreheros}[Extension=.otf,UprightFont=*-regular,BoldFont=*-bold,ItalicFont=*-italic,BoldItalicFont=*-bolditalic]
\setmonofont{texgyrecursor}[Extension=.otf,UprightFont=*-regular,BoldFont=*-bold,Scale=0.88]
\setCJKmainfont{Noto Serif CJK SC}[BoldFont={Noto Serif CJK SC Bold}]
\setCJKsansfont{Noto Sans CJK SC}[BoldFont={Noto Sans CJK SC Bold}]
\setCJKmonofont{Noto Sans Mono CJK SC}
% 带圈数字 ①-⑳ 等划入 CJK 类,用 CJK 字体渲染(xeCJK 默认按拉丁处理)
\xeCJKDeclareCharClass{CJK}{"2460 -> "24FF}

\usepackage{amsmath,amssymb,bm}
\usepackage{booktabs}
\usepackage{graphicx}
\usepackage{xcolor}
\usepackage{titlesec}
\usepackage{fancyhdr}
\usepackage{enumitem}
\usepackage{caption}
\usepackage{listings}

% ---- 配色(主色可改) ----
\definecolor{accent}{HTML}{0E5FA8}     % 强调蓝(链接/节号)
\definecolor{good}{HTML}{0E7A4D}       % 正向结果
\definecolor{warn}{HTML}{B45309}       % 负向/警示
\definecolor{faint}{HTML}{6B7889}      % 弱化文字
\definecolor{codebg}{HTML}{F0F2F5}
\definecolor{rule}{HTML}{C9CFD8}

\usepackage[colorlinks=true,linkcolor=accent,citecolor=accent,urlcolor=accent]{hyperref}

% ---- 章节标题:无衬线粗体,节号着色 ----
\titleformat{\section}{\Large\bfseries\sffamily}{\textcolor{accent}{\thesection}}{0.6em}{}
\titleformat{\subsection}{\large\bfseries\sffamily}{\textcolor{accent}{\thesubsection}}{0.5em}{}
\titleformat{\subsubsection}{\normalsize\bfseries\sffamily}{\textcolor{accent}{\thesubsubsection}}{0.5em}{}
\titlespacing{\section}{0pt}{1.6em}{0.7em}
\titlespacing{\subsection}{0pt}{1.2em}{0.5em}

% ---- 页眉页脚(Cosmos 式:左标题右日期,下细线) ----
\newcommand{\runningtitle}{}  % 主文档用 \renewcommand 设置
\newcommand{\reportdate}{}
\newcommand{\reportfooter}{}
\fancypagestyle{cnreport}{%
  \fancyhf{}
  \fancyhead[L]{\small\sffamily \runningtitle}
  \fancyhead[R]{\small\sffamily \reportdate}
  \fancyfoot[L]{\small\color{faint}\sffamily \reportfooter}
  \fancyfoot[R]{\small\sffamily \thepage}
  \renewcommand{\headrulewidth}{0.4pt}
  \renewcommand{\footrulewidth}{0pt}
}
\pagestyle{cnreport}

% ---- 题注与列表 ----
\captionsetup{font=small,labelfont={bf,sf},skip=6pt}
\setlist{itemsep=2pt,topsep=4pt,parsep=0pt}

% ---- 代码 ----
\lstset{
  basicstyle=\ttfamily\small,
  backgroundcolor=\color{codebg},
  frame=single, framerule=0.3pt, rulecolor=\color{rule},
  xleftmargin=6pt, xrightmargin=6pt,
  breaklines=true, columns=fullflexible, keepspaces=true,
  aboveskip=8pt, belowskip=8pt,
}

% ---- 报告头(kicker / 标题 / 副题 / 署名 / 摘要 / 关键词) ----
% 用法: \reporthead{kicker-left}{kicker-right}{标题}{副题}{署名行}
\newcommand{\reporthead}[5]{%
  \begin{center}
  {\small\sffamily\bfseries\color{accent} #1 \hfill #2}\\[2pt]
  {\color{black}\rule{\linewidth}{1.2pt}}\\[14pt]
  {\LARGE\sffamily\bfseries #3\par}\vspace{8pt}
  {\large\color{faint} #4\par}\vspace{10pt}
  {\small\sffamily\color{faint} #5\par}\vspace{4pt}
  {\color{rule}\rule{\linewidth}{0.4pt}}
  \end{center}\vspace{6pt}
}
% 摘要块: \begin{cnabstract} ... \end{cnabstract}
\newenvironment{cnabstract}{%
  \begin{quote}\small\color{faint}\sffamily\bfseries\color{black} 摘要 \quad\normalfont\rmfamily\color{faint!80!black}%
}{\end{quote}\vspace{2pt}}
% 重点结论框: \begin{quotebox} ... \end{quotebox}
\newenvironment{quotebox}{%
  \begin{quote}\small
  {\color{rule}\rule{\linewidth}{0.8pt}}\par\vspace{2pt}%
}{\par\vspace{2pt}{\color{rule}\rule{\linewidth}{0.8pt}}\end{quote}\vspace{2pt}}

\newcommand{\keywords}[1]{\begin{quote}\small\color{faint}{\bfseries\sffamily\color{black} 关键词:\ } #1\end{quote}}

% 表内彩色标记
\newcommand{\upgood}[1]{\textcolor{good}{\bfseries #1}}
\newcommand{\downbad}[1]{\textcolor{warn}{\bfseries #1}}
\newcommand{\dimtext}[1]{\textcolor{faint}{#1}}

\setlength{\parskip}{2pt}
\linespread{1.12}
\emergencystretch=3em  % CJK 段落容忍度,消除 overfull

% 编译: tectonic(XeTeX 引擎),需要系统安装 Noto Serif/Sans CJK SC 字体
\usepackage[a4paper,top=2.7cm,bottom=2.7cm,left=2.5cm,right=2.5cm,headheight=15pt]{geometry}
\usepackage{fontspec}
% 正文:拉丁衬线 + 思源宋体;标题:拉丁无衬线 + 思源黑体
% 按文件名从 TeX 字体树加载(tectonic bundle 自带 TeX Gyre;系统字体需 Noto CJK SC)
\setmainfont{texgyretermes}[Extension=.otf,UprightFont=*-regular,BoldFont=*-bold,ItalicFont=*-italic,BoldItalicFont=*-bolditalic]
\setsansfont{texgyreheros}[Extension=.otf,UprightFont=*-regular,BoldFont=*-bold,ItalicFont=*-italic,BoldItalicFont=*-bolditalic]
\setmonofont{texgyrecursor}[Extension=.otf,UprightFont=*-regular,BoldFont=*-bold,Scale=0.88]
\setCJKmainfont{Noto Serif CJK SC}[BoldFont={Noto Serif CJK SC Bold}]
\setCJKsansfont{Noto Sans CJK SC}[BoldFont={Noto Sans CJK SC Bold}]
\setCJKmonofont{Noto Sans Mono CJK SC}
% 带圈数字 ①-⑳ 等划入 CJK 类,用 CJK 字体渲染(xeCJK 默认按拉丁处理)
\xeCJKDeclareCharClass{CJK}{"2460 -> "24FF}

\usepackage{amsmath,amssymb,bm}
\usepackage{booktabs}
\usepackage{graphicx}
\usepackage{xcolor}
\usepackage{titlesec}
\usepackage{fancyhdr}
\usepackage{enumitem}
\usepackage{caption}
\usepackage{listings}

% ---- 配色(主色可改) ----
\definecolor{accent}{HTML}{0E5FA8}     % 强调蓝(链接/节号)
\definecolor{good}{HTML}{0E7A4D}       % 正向结果
\definecolor{warn}{HTML}{B45309}       % 负向/警示
\definecolor{faint}{HTML}{6B7889}      % 弱化文字
\definecolor{codebg}{HTML}{F0F2F5}
\definecolor{rule}{HTML}{C9CFD8}

\usepackage[colorlinks=true,linkcolor=accent,citecolor=accent,urlcolor=accent]{hyperref}

% ---- 章节标题:无衬线粗体,节号着色 ----
\titleformat{\section}{\Large\bfseries\sffamily}{\textcolor{accent}{\thesection}}{0.6em}{}
\titleformat{\subsection}{\large\bfseries\sffamily}{\textcolor{accent}{\thesubsection}}{0.5em}{}
\titleformat{\subsubsection}{\normalsize\bfseries\sffamily}{\textcolor{accent}{\thesubsubsection}}{0.5em}{}
\titlespacing{\section}{0pt}{1.6em}{0.7em}
\titlespacing{\subsection}{0pt}{1.2em}{0.5em}

% ---- 页眉页脚(Cosmos 式:左标题右日期,下细线) ----
\newcommand{\runningtitle}{}  % 主文档用 \renewcommand 设置
\newcommand{\reportdate}{}
\newcommand{\reportfooter}{}
\fancypagestyle{cnreport}{%
  \fancyhf{}
  \fancyhead[L]{\small\sffamily \runningtitle}
  \fancyhead[R]{\small\sffamily \reportdate}
  \fancyfoot[L]{\small\color{faint}\sffamily \reportfooter}
  \fancyfoot[R]{\small\sffamily \thepage}
  \renewcommand{\headrulewidth}{0.4pt}
  \renewcommand{\footrulewidth}{0pt}
}
\pagestyle{cnreport}

% ---- 题注与列表 ----
\captionsetup{font=small,labelfont={bf,sf},skip=6pt}
\setlist{itemsep=2pt,topsep=4pt,parsep=0pt}

% ---- 代码 ----
\lstset{
  basicstyle=\ttfamily\small,
  backgroundcolor=\color{codebg},
  frame=single, framerule=0.3pt, rulecolor=\color{rule},
  xleftmargin=6pt, xrightmargin=6pt,
  breaklines=true, columns=fullflexible, keepspaces=true,
  aboveskip=8pt, belowskip=8pt,
}

% ---- 报告头(kicker / 标题 / 副题 / 署名 / 摘要 / 关键词) ----
% 用法: \reporthead{kicker-left}{kicker-right}{标题}{副题}{署名行}
\newcommand{\reporthead}[5]{%
  \begin{center}
  {\small\sffamily\bfseries\color{accent} #1 \hfill #2}\\[2pt]
  {\color{black}\rule{\linewidth}{1.2pt}}\\[14pt]
  {\LARGE\sffamily\bfseries #3\par}\vspace{8pt}
  {\large\color{faint} #4\par}\vspace{10pt}
  {\small\sffamily\color{faint} #5\par}\vspace{4pt}
  {\color{rule}\rule{\linewidth}{0.4pt}}
  \end{center}\vspace{6pt}
}
% 摘要块: \begin{cnabstract} ... \end{cnabstract}
\newenvironment{cnabstract}{%
  \begin{quote}\small\color{faint}\sffamily\bfseries\color{black} 摘要 \quad\normalfont\rmfamily\color{faint!80!black}%
}{\end{quote}\vspace{2pt}}
% 重点结论框: \begin{quotebox} ... \end{quotebox}
\newenvironment{quotebox}{%
  \begin{quote}\small
  {\color{rule}\rule{\linewidth}{0.8pt}}\par\vspace{2pt}%
}{\par\vspace{2pt}{\color{rule}\rule{\linewidth}{0.8pt}}\end{quote}\vspace{2pt}}

\newcommand{\keywords}[1]{\begin{quote}\small\color{faint}{\bfseries\sffamily\color{black} 关键词:\ } #1\end{quote}}

% 表内彩色标记
\newcommand{\upgood}[1]{\textcolor{good}{\bfseries #1}}
\newcommand{\downbad}[1]{\textcolor{warn}{\bfseries #1}}
\newcommand{\dimtext}[1]{\textcolor{faint}{#1}}

\setlength{\parskip}{2pt}
\linespread{1.12}
\emergencystretch=3em  % CJK 段落容忍度,消除 overfull


\renewcommand{\runningtitle}{LingBot-VLA 2.0 LeRobot 移植交接}
\renewcommand{\reportdate}{2026 年 8 月 26 日}
\renewcommand{\reportfooter}{LingBot-VLA 移植工作组 · 内部交接}

\begin{document}

\reporthead{内部交接 · 供代码 Review}{2026 年 8 月 26 日}
{LingBot-VLA 2.0 的 LeRobot 移植:代码逻辑、上游适配与框架适配}
{给内部同学的代码导读——这套移植怎么读、上游权重怎么来的、怎么长在 lerobot 里}
{LingBot-VLA 移植工作组 \quad 代码:\texttt{lingbot-lerobot-final}(私有)主干 \texttt{2902c42} \quad 上游 PR:huggingface/lerobot \#3967}

\begin{cnabstract}
本文是 LingBot-VLA 2.0 LeRobot 移植的代码交接导读,面向要通读并 review 这套代码的内部同学。全文回答三个问题:\textbf{(1) 代码逻辑}——双流架构(Qwen3-VL-4B 主干 + 36 层稀疏 MoE 动作专家)如何组织成 lerobot 的 BYOP 三件套,一次推理的前向三阶段(前缀编码 $\to$ KV 复用去噪 $\to$ 反归一化)各做什么;\textbf{(2) 上游适配}——官方发布的 \texttt{robbyant/lingbot-vla-v2-6b} 权重为何不能直接加载、一次性转换器做了什么、与 transformers 原生实现的逐位数值一致性如何;\textbf{(3) 框架适配}——策略如何零特判地接入 lerobot 的 processor 注册链、训练/推理入口与多卡路线。附一页性能速查与已知边界。数值上与上游逐位等价,CI 全绿,唯一未闭环项为 RoboTwin 闭环 port 级复评。
\end{cnabstract}
\keywords{VLA;LingBot-VLA 2.0;lerobot BYOP;流匹配;稀疏 MoE;checkpoint 转换;逐位一致性}

\section{一页总览}

\textbf{是什么。}LingBot-VLA 2.0 是 pi0 式双流 VLA:Qwen3-VL-4B-Instruct 作视觉-语言主干,并联一个 36 层(hidden 768)的 Qwen2 动作专家,其\textbf{全部 36 层}均为稀疏 MoE(32 专家、top-4 路由、专家中间维 512、共享专家 704);动作由流匹配生成,默认 10 步 Euler 去噪,一次输出 chunk 50。全模型 6.23B 参数。本移植将其以 lerobot BYOP(Build-Your-Own-Policy)形式落为策略 \texttt{lingbot\_vla\_v2}。

\textbf{交付态}(截至 2026-08-26):

\begin{itemize}[nosep,leftmargin=1.4em]
  \item huggingface/lerobot PR \#3967:四轮 review 意见\textbf{全部清零}(细节见仓库内 \texttt{PR3967\_REVIEW\_TODO\_20260826.md}),CI 全绿(pre-commit 19 hooks + pytest 422 passed / 18 skipped);
  \item 私有存档库 \texttt{miracle-techlink/lingbot-lerobot-final}(单提交快照,即本次 review 对象);
  \item 数值一致性:与 transformers 4.57 原生前向\textbf{逐位等价}(唯一残差为 triton fused\_moe kernel 内部 MMA $\approx$0.06,良性)。
\end{itemize}

\textbf{端到端数据流}(一次 \texttt{predict\_action\_chunk}):

\begin{center}\small
\setlength{\fboxsep}{3pt}
\fcolorbox{rule}{codebg}{%
\begin{minipage}{0.94\linewidth}\centering\ttfamily
obs(相机$\times$N / state / task 文本)\\
$\downarrow$~PreProcessor:\texttt{tokenizer} + Qwen 图像预处理(GPU 可选)+ \texttt{FeatureTransform.apply}(slot 映射 + 归一化)\\
$\downarrow$~\texttt{embed\_prefix}:vision tower 编码图像 + 文本 $\to$ 36 层前缀 KV cache(\textbf{每次推理仅一次})\\
$\downarrow$~去噪环 $\times$ \texttt{num\_steps}(默认10):\texttt{predict\_velocity} 逐步走专家流——每层注意力复用前缀 KV + MoE FFN\\
$\downarrow$~速度场 ODE 积分 $\to$ action chunk(1, 50, action\_dim)\\
$\downarrow$~PostProcessor:\texttt{FeatureTransform.unapply}(canonical $\to$ 机器人关节空间,反归一化;缺配置即 raise)
\end{minipage}}
\end{center}

\section{代码逻辑}

\subsection{文件地图}

代码在 \texttt{src/lerobot/policies/lingbot\_vla\_v2/} 下,按 lerobot BYOP 惯例组织:

\begin{table}[htbp]\centering\small
\caption{文件职责一览(review 按此表定位)}
\begin{tabular}{@{}lp{10.6cm}@{}}
\toprule
\textbf{文件} & \textbf{职责} \\
\midrule
\texttt{configuration\_*.py} & 策略配置(draccus 注册,\texttt{type: lingbot\_vla\_v2})。\texttt{resolve\_robot\_config\_and\_stats} 做"显式路径优先、内嵌内容兜底"的资产解析;\texttt{build\_feature\_transform\_configs} 是 apply/unapply 两侧共用的\textbf{单一配置来源};\texttt{\_save\_pretrained} 在保存时清空机器本地路径字段,保证 checkpoint 自包含 \\
\texttt{modeling\_*.py} & \texttt{LingbotVLAV2Policy} 与内部 \texttt{LingBotVLAFlowMatchingModel}:前缀编码、KV cache 生命周期、去噪环、\texttt{\_postprocess\_actions} 反归一化(反归一化器构建失败\textbf{直接 raise},绝不静默输出归一化值——真机上等于任意关节命令) \\
\texttt{processor\_*.py} & ProcessorStep 管线:pre(tokenize + Qwen 图像预处理 + DeviceProcessorStep)/ post(unapply);提供 \texttt{make\_*\_pre\_post\_processors\_from\_pretrained},支持 config 派生的 \texttt{feature\_step\_overrides} 注入 \\
\texttt{feature\_transform.py} & \texttt{FeatureTransform}:55 维 canonical 空间(\texttt{max\_state\_dim}/\texttt{max\_action\_dim} 定长 padding)的 slot 映射、逐 slot 归一化统计、apply/unapply 双向 \\
\texttt{qwen3vl\_in\_vla.py} & vendored Qwen3-VL 塔 + VLM 解码层。改造点:双流 forward 与前缀 KV 注入(这是 transformers 里没有的部分) \\
\texttt{qwen2\_action\_expert.py} & vendored Qwen2 动作专家层:含 fused 布局 MoE 与小 batch 稠密双 GEMM 路径 \\
\texttt{convert\_upstream\_checkpoint.py} & 一次性转换器(见 §3) \\
\texttt{configuration\_*\_internal.py} & 上游兼容配置(\texttt{architectures} 指向 FlowMatching / FlowMatchingV2) \\
\bottomrule
\end{tabular}
\end{table}

\subsection{前向三阶段}

\textbf{(1) 前缀编码}(\texttt{embed\_prefix},每次推理一次):图像经 vision tower 得视觉 token(grid-thw 可由 \texttt{precompute\_grid\_thw} 预计算缓存),与文本 token 拼接,过 VLM 主干各层,把 36 层的 K/V 写入前缀 cache。此后整个去噪过程不再触碰主干权重。

\textbf{(2) 去噪环}(\texttt{predict\_velocity} $\times$ \texttt{num\_steps}):动作侧只有 51 个 suffix token(state 1 + action chunk 50)。每步把带噪动作嵌入送入专家流,每层\textbf{注意力复用前缀 KV}(cross-attention 语义,但实现为同一序列的前缀后缀拼接)、FFN 走 MoE top-4;得到速度场后 Euler 积分。步数是质量-延迟旋钮(4/7/10 实测见 §5)。

\textbf{(3) 反归一化}:输出从 55 维 canonical 空间经 \texttt{unapply} 回到具体机器人关节空间并反归一化。apply 与 unapply 共用 \texttt{build\_feature\_transform\_configs} 生成的同一份配置——两侧若各自手写会漂移(actions 在 A 布局下生成、B 布局下反归一化),这是 review 抓出并修掉的真实风险点。

\subsection{关键机制速记}

\begin{itemize}[nosep,leftmargin=1.4em]
  \item \textbf{MoE 两条路}:checkpoint 原生 fused(按专家维堆叠)布局;当 $B{\times}T \le 512$ 时自动走稠密双 GEMM 纯 PyTorch 路径——两个 GEMM 算完 32 专家、top-4 权重折叠进 down GEMM,无路由 kernel、静态形状、dynamo 可整图追踪(与精确路由代数等价,fp32 相对差 1.4e-6)。
  \item \textbf{attention 三实现}:eager / sdpa / flex(flex$\to$triton 编译)。训练推荐 sdpa+bf16;推理 compile 下 eager 与 sdpa 等价;\texttt{fa2} 在上游即 NotImplementedError。
  \item \textbf{加速旋钮}(默认关):\texttt{use\_cudagraph\_denoise}(CUDA graph 捕去噪环,空 cache 捕获)、\texttt{preprocess\_device=cuda}、\texttt{compile\_predict\_velocity}。终烘配方见 §5。
\end{itemize}

\section{上游适配:robbyant 权重 $\to$ LeRobot checkpoint}

\textbf{为什么必须转换。}官方发布物 \texttt{robbyant/lingbot-vla-v2-6b} 的 \texttt{config.json} 只有一行 \texttt{\{"vlm\_family": "qwen3\_v1"\}}——全部架构信息住在官方训练脚本里,任何标准加载器都无法重建模型。早期版本用过"运行时 sniffer + raw loader",但 reviewer 正确指出这会把 4 个特判分支打进 lerobot 核心文件,且 sniffer 会误认别人的 checkpoint;最终方案是\textbf{一次性转换 + 零特判}。

\textbf{转换器做什么}(\texttt{convert\_upstream\_checkpoint.py},跑一次):
\begin{itemize}[nosep,leftmargin=1.4em]
  \item 重打 \texttt{config.json}:\texttt{type: lingbot\_vla\_v2} + 全部结构 override——从此走标准 \texttt{from\_pretrained},多卡加载自动带 \#4042 修复;
  \item 权重切分:1708 张量 $\to$ \textbf{1624 加载 + 84 显式剔除}(depth 对齐 / 视频蒸馏头,\texttt{use\_depth=False};在转换期一次剔除并记录,而非每次加载静默跳过);strict 加载 0 missing / 0 unexpected;
  \item 内嵌机器人资产:\texttt{robot\_config} 与 \texttt{norm\_stats} 内容写进 config 与 processor JSON,保存时清空路径字段 $\Rightarrow$ checkpoint 拷到任何机器都能重载(自包含)。
\end{itemize}

\textbf{数值一致性}(转换正确性的证据):同权重同输入同噪声,transformers 4.57 原生实现 vs 本移植,前缀 embeddings、全层 KV cache、suffix/时间嵌入、MoE router logits、top-4 选中专家全部 \textbf{max$|\Delta|$=0}(bf16 与 fp32 逐位一致)。唯一残差在 vendored triton \texttt{fused\_moe} 的 SwiGLU 内部($\approx$0.06/块),属 kernel 内部低精度 MMA 的良性差异,非结构或加载缺陷。

\textbf{与 transformers 的复用边界}:能用 transformers 的(Qwen3-VL 配置与构件、Qwen2 层、tokenizer/processor)直接用;留在 LeRobot 的是上游没有的部分——双流 forward、前缀 KV 注入、流匹配动作头、MoE 专家逻辑。vendored 文件顶部有来源标注。

\section{LeRobot 框架适配}

\textbf{BYOP 合规}:\texttt{policies/lingbot\_vla\_v2/} 三件套(configuration / modeling / processor)+ \texttt{README.md} 软链到 \texttt{docs/source/policy\_lingbot\_vla\_v2\_README.md} + \texttt{lingbot\_vla\_v2.mdx}(含 Benchmark 章节与复现命令)+ 根 README 模型表 + modelcard 模板条目。

\textbf{Processor 体系接入}:
\begin{itemize}[nosep,leftmargin=1.4em]
  \item \texttt{ProcessorStep} 注册走\textbf{包导入链}:包内 \texttt{\_\_init\_\_.py} 导出 configuration/modeling/processor 三模块(与其他 21 个策略一致),\texttt{import lerobot.policies} 即注册全部 60 个 step,含 \texttt{lingbot\_vla\_2\_feature\_transform}——factory 里为此曾存在的通用 importlib 补丁块已删除;
  \item 管线经 \texttt{PolicyProcessorPipeline.from\_pretrained} 加载;config 派生的 \texttt{feature\_step\_overrides} 注入,使\textbf{新本体 fine-tune 能覆盖 checkpoint 里的槽映射};
  \item lerobot 通用 normalizer override 在此被过滤——本策略归一化在 \texttt{FeatureTransform} 内,不另设 normalizer step(过滤逻辑在 lingbot processor 模块内,不污染共享代码)。
\end{itemize}

\textbf{零特判}:\texttt{configs/policies.py}、\texttt{factory.py}、\texttt{lerobot\_train.py}、\texttt{lerobot\_eval.py} 四个核心文件中 lingbot 特判为 \textbf{0 处};processor 工厂按 isinstance 分发(与 Groot/Evo1 同款);\texttt{\_no\_split\_modules} 登记真实层类名,FSDP \texttt{TRANSFORMER\_BASED\_WRAP} 开箱即用。

\textbf{训练 / 推理入口}:
\begin{table}[htbp]\centering\small
\begin{tabular}{@{}llp{7.2cm}@{}}
\toprule
\textbf{入口} & \textbf{路线} & \textbf{状态} \\
\midrule
\texttt{lerobot-train} & FSDP1 & 8$\times$A100 干净起训(5.8B/6.2B 可训练,loss $\sim$0.6) \\
 & FSDP2 & 2$\times$4090 冒烟通过(auto-wrap + offload) \\
 & DDP & 2$\to$4 卡线性扩展,8 卡 30k 步 $\approx$ 9--11 h \\
 & LoRA & 社区 PEFT 三钩子路线,单卡 12.4 GB,冒烟通过 \\
\texttt{lerobot-eval} / \texttt{lerobot-rollout} & 标准 & \texttt{predict\_action\_chunk} 返回 (1, 50, 6),数据集单位 \\
\bottomrule
\end{tabular}
\end{table}

\section{性能速查与已知边界}

\begin{table}[htbp]\centering\small
\caption{性能速查(详报:\texttt{lingbot\_vla\_v2\_total\_report.pdf} 与 \texttt{$\sim$/reports/} 汇总)}
\begin{tabular}{@{}lll@{}}
\toprule
\textbf{项} & \textbf{数字} & \textbf{对照} \\
\midrule
数值一致性 & 1624 keys 0/0;前缀/KV/router 全逐位 & transformers 4.57 原生 \\
训练 e2e(A100, B=2) & 2378 $\to$ \upgood{1137 ms(-52\%)} & PR 原版 $\to$ sdpa+bf16+fused AdamW \\
训练吞吐 & B=3 \upgood{1.99 samples/s};LoRA 12.4 GB/卡 & A100-80G \\
推理 e2e(A100) & 1705 $\to$ \upgood{436.6 ms(3.9$\times$)} & 官方 deploy eager $\to$ 本移植优化栈 \\
推理(4090, 终烘) & \upgood{130.17 ms}(7步+CUDA graph) & 官方 130 ms 口径复现;4/7/10 步 = 108.5/130.9/153.0 \\
推理(GB10) & 635 / 752 ms(4步/7步) & 带宽地板:273 GB/s,130ms 物理不可达 \\
真机 RTC 云端推理 & 60 s 49 次 0 错;RTT p50 822 $\to$ 630 ms & q75 量化 + threshold 30 + horizon 24 \\
\bottomrule
\end{tabular}
\end{table}

\textbf{已知边界}(review 时请知悉):
\begin{itemize}[nosep,leftmargin=1.4em]
  \item \textbf{RoboTwin 闭环 port 级复评未完成}:官方权重上游成绩 Clean 93.52\% / Randomized 92.80\%;本移植的闭环初评 0/12 已定位为评测基建问题(评测 ckpt 内嵌慢配置),待快配置重导出后重评——这是当前\textbf{唯一}未闭环的验证项;
  \item depth / 视频蒸馏分支在本 checkpoint 不可达(\texttt{use\_depth=False},84 张量已显式剔除),如需启用走 follow-up;
  \item reviewer 认可延后的 8 项 cosmetic follow-up(process-wide transformers patch、asserts 收敛、license 头、文件重组等)见 TODO 文档 queued 一节;
  \item 四轮 PR review 已全部清零,逐条对照表:\texttt{PR3967\_REVIEW\_TODO\_20260826.md}(仓库根)。
\end{itemize}

\section*{附录:代码位置速查}
\begin{table}[htbp]\centering\small
\begin{tabular}{@{}p{4.6cm}p{9.4cm}@{}}
\toprule
\textbf{想看什么} & \textbf{去哪} \\
\midrule
双流前向 / KV 注入 & \texttt{modeling\_lingbot\_vla\_v2.py}(\texttt{LingBotVLAFlowMatchingModel})+ \texttt{qwen3vl\_in\_vla.py} \\
去噪环 / 速度场 & \texttt{modeling} 中 \texttt{predict\_velocity};CUDA graph 路径同文件 \\
图像 / 文本预处理 & \texttt{processor\_lingbot\_vla\_v2.py}(管线)+ \texttt{feature\_transform.py}(apply) \\
归一化 / slot 映射 & \texttt{feature\_transform.py}(\texttt{FeatureTransform.apply/unapply}) \\
MoE 专家实现 & \texttt{qwen2\_action\_expert.py}(fused 布局 + 稠密双 GEMM) \\
checkpoint 转换 & \texttt{convert\_upstream\_checkpoint.py} \\
训练配方 / 4090 配置 & \texttt{docs/source/lingbot\_vla\_v2.mdx};\texttt{bench/robotwin/} \\
\midrule
外部指针 & 私有库 \texttt{lingbot-lerobot-final};PR \#3967;\texttt{$\sim$/reports/lingbot-vla-v2-训练推理评测汇总-20260826.md};加速详报 \texttt{docs/accel/lingbot\_vla\_v2\_total\_report.pdf} \\
\bottomrule
\end{tabular}
\end{table}

\end{document}
